\documentclass[10pt,a4paper]{article}
% ---------- common preamble : CM3191 organic documents (XeLaTeX) ----------
\usepackage[a4paper,top=16mm,bottom=16mm,left=16mm,right=16mm,headheight=14pt]{geometry}
\usepackage{fontspec}
\usepackage{xeCJK}
\setmainfont{TeX Gyre Termes}
\setsansfont{TeX Gyre Heros}
\setCJKmainfont{Noto Serif CJK KR}
\setCJKsansfont{Noto Sans CJK KR}
\xeCJKsetup{CJKspace=true, CJKecglue={}}
\usepackage{amsmath,amssymb}
\usepackage[version=4]{mhchem}
\usepackage{chemfig}
\usepackage{siunitx}
\sisetup{detect-all, per-mode=symbol, range-phrase=--, range-units=single}
\usepackage{tikz}
\usetikzlibrary{arrows.meta,calc,positioning,shapes.geometric,decorations.pathmorphing,patterns}
\usepackage{tabularx,array,multirow,booktabs}
\usepackage{enumitem}
\usepackage{titlesec}
\usepackage{fancyhdr}
\usepackage{xcolor}
\usepackage{pifont}
\usepackage[hidelinks]{hyperref}
\setlength{\parindent}{0pt}
\setlength{\parskip}{3pt}
\linespread{1.08}

\newcommand{\doctag}{}
\pagestyle{fancy}\fancyhf{}
\lhead{\footnotesize CM3191 Organic \textemdash\ \doctag}
\rhead{\footnotesize KIM DONGHWAN \textbullet\ B7 \textbullet\ Thu}
\cfoot{\footnotesize \thepage}
\renewcommand{\headrulewidth}{0.4pt}

\titleformat{\section}{\large\bfseries}{\thesection}{0.6em}{}
\titlespacing*{\section}{0pt}{9pt}{4pt}
\titleformat{\subsection}{\normalsize\bfseries}{\thesubsection}{0.5em}{}
\titlespacing*{\subsection}{0pt}{7pt}{3pt}

% ---------- form primitives ----------
\newcommand{\wl}[1][4cm]{\rule[-0.6ex]{#1}{0.4pt}}
\newcommand{\cbox}{$\square$}
\newcommand{\tr}{\rule[-3mm]{0pt}{9mm}}
\newcommand{\trs}{\rule[-2mm]{0pt}{7mm}}
\newcommand{\ruled}[1]{%
  \par\vspace{3mm}%
  \foreach \i in {1,...,#1}{\noindent\rule{\linewidth}{0.35pt}\par\vspace{7mm}}%
  \vspace{-7mm}}
\newcommand{\confirm}[1]{\textbf{[CONFIRM: #1]}}

% ---------- boxed callout ----------
\newenvironment{callout}[1]
  {\par\vspace{2mm}\noindent\begin{tabular}{|@{\hspace{2.5mm}}p{\dimexpr\linewidth-7.4mm\relax}@{\hspace{2.5mm}}|}\hline
   \rule{0pt}{4.5mm}\textbf{#1}\par\vspace{0.5mm}\small}
  {\par\vspace{1.5mm}\\\hline\end{tabular}\par\vspace{2mm}}

% ---------- title block ----------
\newcommand{\manualtitle}[3]{%
\begin{center}
{\Large\bfseries #1}\\[1.5mm]
{\large #2}\\[2mm]
{\footnotesize #3}
\end{center}
\vspace{-1mm}\rule{\linewidth}{1pt}\vspace{1mm}
}

% ---------- GTA sign-off ----------
\newcommand{\gtasign}{%
\par\vspace{3mm}
\renewcommand{\arraystretch}{1.4}
\noindent\begin{tabularx}{\linewidth}{|X|X|X|}
\hline
\multicolumn{3}{|l|}{\textbf{To be completed at the end of the session}}\\ \hline
Student signature & GTA name & GTA signature \\ \hline
\tr & & \\ \hline
\end{tabularx}
\renewcommand{\arraystretch}{1}
\par\vspace{1mm}
{\footnotesize This is a working log. The \textbf{official} datasheet (lab manual p.\,7) must still be printed, signed by the GTA and submitted with the report.}
}

% ---------- chemfig defaults ----------
\setchemfig{atom sep=1.9em, bond style={line width=0.6pt}}

\renewcommand{\doctag}{Expt.\ 1 \textemdash\ Concepts}
\setlength{\parskip}{4pt}

\begin{document}

\manualtitle{ORGANIC EXPERIMENT 1 \textemdash\ 개념 정리}
{물에서의 Suzuki--Miyaura 교차 커플링: 원리, 계산, 분석 해석}
{근거: CM3191 Organic Lab Manual AY26/27 (Exp.\,1, pp.\,2--7), CM3191 Org L1 슬라이드, 그리고 본문 끝 참고문헌. 매뉴얼에 없는 수치는 직접 계산했고 계산 과정을 함께 적었습니다.}

\begin{callout}{한 장 요약}
\begin{itemize}[leftmargin=4mm,itemsep=0.5mm,topsep=0.5mm]
\item \textbf{반응}: 4-bromobenzoic acid + phenylboronic acid \textrightarrow\ 4-phenylbenzoic acid(biphenyl-4-carboxylic acid). Pd 촉매가 두 아릴 조각을 Pd 위에 모은 뒤 새 C--C 단일결합으로 연결합니다(L1 p.\,4).
\item \textbf{한계 시약}: 4-bromobenzoic acid 1.69 mmol. \ce{PhB(OH)2} 1.21 eq, \ce{Na2CO3} 2.96 eq, Pd 0.0148 mol\% (= 148 ppm). \textbf{이론 수득량 0.335 g}.
\item \textbf{물에서 되는 이유}: 염기 조건에서 두 기질이 모두 음이온(카복실산염, 붕산염)이 되어 물에 녹고, 촉매 리간드도 수용성입니다. 반응 후에는 산을 넣어 pH를 바꾸는 것만으로 생성물이 침전됩니다.
\item \textbf{분석 포인트}: IR의 카복실산 O--H(넓음)와 C=O. \textsuperscript{1}H NMR에서 치환되지 않은 페닐기의 5H 패턴(특히 1H 삼중선)이 출발물질과 생성물을 구별하는 핵심입니다.
\end{itemize}
\end{callout}

% =====================================================================
\section{촉매 순환 (catalytic cycle)}

\begin{center}
\begin{tikzpicture}[font=\small, >={Stealth[length=2.6mm]}, line width=0.7pt]
\node[draw, rounded corners, inner sep=4pt, fill=white] (pd0) at (0,2.8) {\ce{Pd^0L_n}};
\node[draw, rounded corners, inner sep=4pt, fill=white] (oa)  at (4.6,0) {\ce{Ar-Pd^{II}L_n-Br}};
\node[draw, rounded corners, inner sep=4pt, fill=white] (oh)  at (0,-2.8) {\ce{Ar-Pd^{II}L_n-OH}};
\node[draw, rounded corners, inner sep=4pt, fill=white] (tm)  at (-4.6,0) {\ce{Ar-Pd^{II}L_n-Ph}};
\draw[->] (pd0) to[bend left=28] (oa);
\draw[->] (oa) to[bend left=28] (oh);
\draw[->] (oh) to[bend left=28] (tm);
\draw[->] (tm) to[bend left=28] (pd0);
% step labels (inside)
\node[align=center, font=\footnotesize] at (2.1,1.0) {\textbf{\ding{172} 산화적 첨가}\\oxidative addition\\Pd(0)\,\textrightarrow\,Pd(II)};
\node[align=center, font=\footnotesize] at (2.1,-1.0) {\textbf{\ding{173} 리간드 교환}\\\ce{Br^-} \textrightarrow\ \ce{OH^-}};
\node[align=center, font=\footnotesize] at (-2.1,-1.0) {\textbf{\ding{174} 금속교환}\\transmetalation\\Ph: B\,\textrightarrow\,Pd};
\node[align=center, font=\footnotesize] at (-2.1,1.0) {\textbf{\ding{175} 환원적 제거}\\reductive elimination\\Pd(II)\,\textrightarrow\,Pd(0)};
% inputs / outputs (outside)
\node (arbr) at (5.6,2.7) {\ce{Ar-Br} in};
\draw[->, dashed] (arbr.west) to[bend right=10] (3.05,2.0);
\node (prod) at (-5.6,2.7) {\ce{Ar-Ph} out};
\draw[->, dashed] (-3.05,2.0) to[bend left=10] (prod.east);
\node (base) at (5.8,-2.6) {\ce{OH^-} in};
\draw[->, dashed] (base.west) to[bend left=10] (3.05,-2.0);
\node (brm) at (3.4,-3.7) {\ce{Br^-} out};
\draw[->, dashed] (2.3,-2.45) to[bend right=10] (brm.north west);
\node[align=center] (boron) at (-6.1,-2.6) {\ce{PhB(OH)2} in\\{\footnotesize (또는 \ce{PhB(OH)3^-})}};
\draw[->, dashed] (boron.east) to[bend right=10] (-3.05,-2.0);
\node (boh3) at (-3.4,-3.7) {\ce{B(OH)3} out};
\draw[->, dashed] (-2.3,-2.45) to[bend left=10] (boh3.north east);
\end{tikzpicture}\\[1mm]
{\footnotesize Ar = \ce{4-(^{-}O2C)C6H4}; \ L = 탈양성자화된 2-amino-4,6-dihydroxypyrimidine 리간드; \ Pd 0.0148 mol\%. 금속교환은 oxo-palladium 경로로 그렸습니다(본문 참조).}
\end{center}

\textbf{\ding{172} 산화적 첨가.} Pd(0)이 \ce{C-Br} 결합 사이로 끼어들어 Pd의 산화수가 2 늘고 리간드 두 개(Ar, Br)를 얻습니다(L1 p.\,5). 아릴 할라이드의 반응성은 \textbf{Ar--I $>$ Ar--Br $\approx$ Ar--OTf $\gg$ Ar--Cl $\gg$ Ar--F}입니다(L1 p.\,10). 이 순서는 C--X 결합이 강해지는 순서와 반대입니다. 결합이 강할수록 산화적 첨가가 느려지기 때문입니다. 아릴 브로마이드는 반응성과 가격 사이의 절충점이라 교육용 실험에서 흔히 씁니다. 반면 아릴 클로라이드는 보통 전자가 풍부하고 부피가 큰 phosphine(예: XPhos)이 필요합니다. L1 연습문제(pimodivir)에서 클로라이드 \textbf{14}의 커플링에 \ce{Pd(OAc)2}/X-Phos를 쓴 것도 같은 맥락입니다(8절).

\textbf{\ding{173} \& \ding{174} 금속교환과 염기.} 유기붕소 화합물은 공유결합성이 커서 그대로는 금속교환이 잘 일어나지 않습니다. 그래서 염기가 반드시 필요합니다(L1 p.\,9). 이 단계에는 두 가지 경로가 제안되어 있습니다.
\begin{itemize}[leftmargin=5mm,itemsep=0.5mm,topsep=0.5mm]
\item \textbf{Boronate 경로 (L1의 설명)}: \ce{PhB(OH)2 + OH^- -> PhB(OH)3^-}. 음전하를 띤 4배위 ``ate'' 착물은 Ph기의 친핵성이 커서 \ce{Ar-Pd-Br}의 Br을 치환하며 Ph를 넘겨줍니다.
\item \textbf{Oxo-palladium 경로}: 먼저 \ce{Ar-Pd-Br}이 \ce{Ar-Pd-OH}로 바뀌고, 이 Pd--OH가 중성 보론산과 Pd--O--B 다리를 만든 뒤 Ph가 옮겨갑니다(위 그림은 이 경로로 그렸습니다).
\end{itemize}
Carrow와 Hartwig(2011)은 분리한 Pd 착물로 두 반응의 속도를 비교했습니다. 그 결과 \textbf{약염기·수계 조건에서는 Pd--OH + 보론산 경로가 우세하다}고 결론지었습니다. 반면 일부 계산 연구는 boronate 경로의 장벽이 더 낮다고 보고해, 조건에 따라 달라지는 미해결 쟁점으로 남아 있습니다. \textbf{시험이나 proforma에서는 강의 설명(boronate)을 기본으로 쓰고}, 고찰에서 다른 경로를 언급하는 정도가 적절해 보입니다.

\textbf{\ding{175} 환원적 제거.} Pd 위에서 cis로 놓인 Ar과 Ph가 결합하며 떨어져 나오고, Pd(0)이 재생됩니다. 새 σ결합은 두 ``옛 리간드'' 사이에 생깁니다(L1 p.\,5).

\textbf{전촉매 활성화.} 투입한 것은 \ce{Pd(OAc)2}, 즉 Pd(II)입니다. 순환에 들어가려면 먼저 Pd(0)으로 환원되어야 합니다. phosphine이 없는 계에서는 보론산 두 분자가 차례로 금속교환한 뒤 환원적 제거로 \ce{Ph-Ph}(biphenyl, 호모커플링)를 내놓으며 Pd(0)이 생기는 경로가 흔히 제안됩니다. 다만 Pd 양이 0.25 µmol에 불과해 이로 인한 biphenyl도 최대 0.25 µmol 수준이며, 수율에는 사실상 영향이 없습니다.

\subsection*{전체 반응식 (염기 조건, 물질수지 확인)}
\[
\ce{Br-C6H4-CO2^- + PhB(OH)2 + OH^- -> Ph-C6H4-CO2^- + B(OH)3 + Br^-}
\]
{\small 좌변과 우변의 원자(C, H, B, Br, O)와 전하($-2$)가 일치합니다. 염기 조건이므로 \ce{B(OH)3}는 다시 \ce{B(OH)4^-}로 존재합니다. \ce{OH^-}는 탄산염의 가수분해(\ce{CO3^2- + H2O <=> HCO3^- + OH^-})에서 공급됩니다.}

% =====================================================================
\section{이 실험에서 염기(\ce{Na2CO3})가 하는 일}

\ce{Na2CO3}는 5.0 mmol로 2.96 eq입니다. 촉매 순환의 염기 역할 외에도 세 가지 일을 동시에 합니다.
\begin{enumerate}[leftmargin=6mm,itemsep=0.5mm]
\item \textbf{가용화}: 4-bromobenzoic acid(p$K_\text{a}$는 benzoic acid의 4.2와 비슷)를 카복실산염으로 바꿔 물에 녹입니다. 매뉴얼이 ``모든 고체가 녹을 때까지'' 촉매 투입을 미루는 이유입니다.
\item \textbf{붕소 활성화}: 보론산을 붕산염(\ce{PhB(OH)3^-})으로 만들거나, Pd--OH를 형성하게 합니다.
\item \textbf{부산물 중화}: 반응에서 생기는 HBr 등가물과 \ce{B(OH)3}를 중화해 pH를 유지합니다.
\end{enumerate}
\textbf{초기 pH 추정}: \ce{CO3^2-}의 $K_\text{b} = K_\text{w}/K_\text{a2} = 10^{-14}/10^{-10.33} = 2.1\times10^{-4}$입니다. 따라서 $[\ce{OH^-}] \approx \sqrt{K_\text{b}\times0.5} = 1.0\times10^{-2}$ M, 즉 \textbf{pH $\approx$ 12}입니다. 카복실산 중화로 일부 소모되어도 반응 중 강한 염기성이 유지됩니다.

% =====================================================================
\section{왜 물인가 (L1 p.\,10의 두 질문)}

\textbf{Q1. 물을 용매로 쓸 때의 장점과 어려움은?}\par
\renewcommand{\arraystretch}{1.2}
\begin{tabularx}{\linewidth}{|X|X|}
\hline
\textbf{장점} & \textbf{어려움} \\ \hline
\footnotesize 무독성·불연성·저렴. toluene, DMF 같은 기존 용매를 대체합니다(매뉴얼 p.\,2). \newline
유기붕소 부산물(붕산염)과 무기염이 수용액에 남아 분리가 쉽습니다(L1 p.\,7). \newline
공기와 물에 안정한 보론산, phosphine이 없는 촉매 덕분에 불활성 분위기가 필요 없습니다. \newline
pH만 바꿔 생성물을 침전시킬 수 있어 추출 용매가 필요 없습니다.
&
\footnotesize 대부분의 유기 기질은 물에 녹지 않습니다. \newline
전형적인 Pd--phosphine 촉매는 물에 녹지 않으므로 수용성 리간드가 필요합니다. \newline
고온·강염기 수용액에서는 protodeboronation(\ce{PhB(OH)2 -> PhH})이 경쟁할 수 있습니다. \newline
물은 증발 잠열이 커서 제거에 에너지가 많이 듭니다. \newline
Pd가 수계 폐액으로 들어가므로 금속 폐액으로 따로 처리해야 합니다.
\\ \hline
\end{tabularx}
\renewcommand{\arraystretch}{1}

\textbf{Q2. 이 반응은 왜 물에서 잘 되는가?}
\begin{itemize}[leftmargin=5mm,itemsep=0.5mm,topsep=0.5mm]
\item \textbf{기질 설계}: 할로젠화물에 COOH라는 ``수용성 손잡이''가 있습니다. 염기에서 카복실산염이 되고, 보론산도 붕산염이 되어 \textbf{두 반응물 모두 균일 수용액에} 존재합니다.
\item \textbf{촉매 설계}: 2-amino-4,6-dihydroxypyrimidine은 NaOH로 탈양성자화되면(슬라이드 11의 NaO-치환 구조) 물에 녹습니다. 이 리간드로 만든 Pd 착물은 phosphine 없이도 수용액에서 Suzuki 커플링을 촉매한다고 보고되었습니다(Chalker, Wood, Davis, 2009). 이 실험은 그 촉매를 교육용으로 옮긴 Hamilton 등(2013)의 절차를 따릅니다.
\item \textbf{분리 설계}: 생성물도 염기에서는 녹아 있다가, 산을 넣으면 중성 카복실산이 되어 찬물에서 침전합니다. ``녹여서 반응시키고, pH로 꺼낸다''는 구조입니다.
\end{itemize}

% =====================================================================
\section{촉매량, TON, TOF}

\textbf{촉매량}: Pd $= 1.0\ \text{mL}\times0.25\ \text{mM} = 0.25\ \mu\text{mol}$입니다. 기질 1.691 mmol 대비 \textbf{0.0148 mol\%}로, 슬라이드의 ``0.015\%''와 일치합니다. 몰비로는 148 ppm입니다.

\textbf{정의 (L1 p.\,19)}:
\[
\text{TON} = \frac{n_\text{product}}{n_\text{Pd}}, \qquad
\text{TOF} = \frac{\text{TON}}{t}
\]
TON은 무차원이고, TOF 단위는 \si{\per\second}, \si{\per\minute}, \si{\per\hour} 중 하나로 명시해야 합니다.

\begin{center}
\renewcommand{\arraystretch}{1.2}
\begin{tabular}{|c|c|c|c|c|}
\hline
\textbf{분리 수율} & \textbf{생성물 (g)} & \textbf{TON} & \textbf{TOF (\si{\per\minute}, 30 min)} & \textbf{TOF (\si{\per\hour})} \\ \hline
50\% & 0.168 & 3383 & 113 & 6765 \\ \hline
70\% & 0.235 & 4736 & 158 & 9472 \\ \hline
90\% & 0.302 & 6089 & 203 & 12178 \\ \hline
100\% (상한) & 0.335 & 6765 & 226 & 13531 \\ \hline
\end{tabular}
\renewcommand{\arraystretch}{1}
\end{center}
{\small \textbf{보고서에서 밝혀야 할 가정.} (i) \textbf{분리 수율 기준}이므로 재결정 손실만큼 TON이 과소평가됩니다. 전환율 기준 TON은 이보다 큽니다. (ii) 투입한 Pd가 \textbf{전부 활성}이라고 가정했습니다. 실제 활성종은 일부일 수 있어 ``site당'' 값은 더 클 수 있습니다. (iii) 반응이 30분보다 일찍 끝났다면 TOF는 과소평가됩니다. 즉 이 값은 30분 평균값입니다.}

% =====================================================================
\section{후처리: 산-염기 화학으로 보는 각 단계}

\textbf{산성화 (1 M HCl 20 mL = 20 mmol).}
\[
\ce{CO3^2- + 2 H+ -> CO2 ^ + H2O}, \qquad \ce{Ph-C6H4-CO2^- + H+ -> Ph-C6H4-CO2H v}
\]
\begin{itemize}[leftmargin=5mm,itemsep=0.5mm,topsep=0.5mm]
\item 탄산염 5.0 mmol을 \ce{CO2}까지 없애는 데 필요한 산은 10 mmol입니다. HCl 20 mmol은 이보다 약 10 mmol 많습니다. 반응 중 카복실산이 내놓은 양성자는 산성화 때 도로 붙으므로 알짜 효과가 없습니다.
\item 과량 \ce{H+} 약 10 mmol이 부피 약 51 mL(반응액 11 + HCl 20 + 물 20)에 녹아 있으므로 $[\ce{H+}]\approx 0.20$ M, \textbf{pH $\approx$ 0.7}로 계산됩니다. 생성물(p$K_\text{a}\approx4$)은 사실상 100\% 양성자화됩니다.
\item \textbf{기체}는 \ce{CO2}입니다. HCl을 한 방울씩 넣는 이유는 거품이 넘치는 것을 막기 위해서입니다. \textbf{얼음 수조}는 중화열을 흡수하고 생성물의 용해도를 낮춥니다.
\item 물 20 mL를 넣고 10분 교반하면 NaCl, 붕산, 남은 보론산 같은 수용성 물질이 녹아 나가고, 침전이 숙성되어 거르기 쉬워집니다.
\end{itemize}

\textbf{재결정 (에탄올, HCl 3 mL 첨가).}
\begin{itemize}[leftmargin=5mm,itemsep=0.5mm,topsep=0.5mm]
\item 원리는 \textbf{뜨거울 때는 잘 녹고, 차가울 때는 덜 녹는} 용매 하나를 쓰는 것입니다. 가장 흔한 수율 손실 원인은 뜨거운 에탄올을 너무 많이 넣는 것입니다.
\item \textbf{HCl을 넣는 이유}: 조생성물에 남은 카복실산염이나 탄산염을 중성 산으로 바꿔 한 가지 형태로 결정화시키기 위해서입니다. 염 형태로 남으면 용해도와 녹는 거동이 달라집니다.
\item \textbf{불순물의 행방}: biphenyl(호모커플링)은 에탄올에 잘 녹아 모액에 남습니다. 붕소 부산물도 대부분 모액이나 세척액으로 빠집니다. 반면 \textbf{미반응 4-bromobenzoic acid는 생성물과 성질이 비슷해 일부 함께 결정화될 수 있으므로} NMR로 확인해야 합니다.
\item \textbf{천천히 식히는 이유}: 결정이 천천히 자라야 격자에 불순물이 덜 끼어 순도가 높아집니다.
\item \textbf{수율이 100\%를 넘는다면} 건조가 덜 되어 물이나 에탄올이 남았거나 무기염이 섞였다는 뜻입니다. IR의 넓은 O--H와 NMR의 \ce{H2O}/EtOH 신호로 뒷받침할 수 있습니다.
\end{itemize}

% =====================================================================
\section{예상 분석 결과와 해석}

\subsection*{IR (ATR)}
아래는 일반적인 작용기 흡수 범위이며(Pavia 등), 실제 값은 본인 스펙트럼에서 읽어야 합니다.
\begin{center}
\renewcommand{\arraystretch}{1.15}
\begin{tabularx}{\linewidth}{|l|c|X|}
\hline
\textbf{진동} & \textbf{범위 (\si{\per\centi\metre})} & \textbf{해석} \\ \hline
O--H 신축 (카복실산 이량체) & 3300--2500, 매우 넓음 & 수소결합 이량체의 특징. C--H 신축과 겹칩니다 \\ \hline
C=O 신축 & 약 1710--1680 & 방향족 고리와 짝지어져 포화 카복실산보다 낮은 쪽에 나옵니다 \\ \hline
방향족 C=C & 약 1600, 1500--1450 & 두 벤젠 고리 \\ \hline
C--O 신축 & 1320--1210 & 카복실산 \\ \hline
C--H 면외 굽힘 & 약 860--800 / 770--730 \& 710--690 & para-이치환 고리 / 일치환 페닐 고리 \\ \hline
\end{tabularx}
\renewcommand{\arraystretch}{1}
\end{center}
{\small 출발물질에도 카복실산 신호가 있으므로 IR만으로 커플링 여부를 증명하기는 어렵습니다. \textbf{일치환 페닐의 면외 굽힘(약 750, 690)}이 새로 나타나는지가 보조 근거가 됩니다.}

\subsection*{\textsuperscript{1}H NMR (DMSO-$d_6$, 문헌 피크 목록 기준)}
\begin{center}
\renewcommand{\arraystretch}{1.15}
\begin{tabularx}{\linewidth}{|c|c|c|X|}
\hline
\textbf{δ (ppm)} & \textbf{다중도} & \textbf{적분} & \textbf{귀속} \\ \hline
13.05 & br s & 1H & COOH (교환성, 넓거나 약하게 보일 수 있음) \\ \hline
8.07 & d, $J\approx8.4$ Hz & 2H & 가운데 고리, COOH의 ortho H (AA$'$BB$'$) \\ \hline
7.82 & d, $J\approx8.4$ Hz & 2H & 가운데 고리, 페닐 쪽 H (AA$'$BB$'$) \\ \hline
7.75 & d & 2H & 말단 페닐 ortho H \\ \hline
7.52 & t & 2H & 말단 페닐 meta H \\ \hline
7.44 & t & 1H & 말단 페닐 para H \\ \hline
\end{tabularx}
\renewcommand{\arraystretch}{1}
\end{center}
{\small 출처: ChemicalBook 스펙트럼 DB, 4-Biphenylcarboxylic acid (CAS 92-92-2), 400 MHz 피크 목록에서 봉우리 중심을 읽은 값. 목록에 용매 표기는 없지만 COOH가 13.05 ppm에 관측되므로 DMSO-$d_6$로 판단됩니다. 본인 스펙트럼은 농도·온도에 따라 약간 다를 수 있으므로 보고서에는 본인 값을 쓰고 이 값과 비교하세요. 가운데 고리의 para-이치환 패턴은 엄밀히 AA$'$BB$'$ 스핀계이지만 겉보기에는 두 개의 ``doublet''으로 보입니다.}

\textbf{출발물질과 불순물 판별.}
\begin{itemize}[leftmargin=5mm,itemsep=0.5mm,topsep=0.5mm]
\item \textbf{4-Bromobenzoic acid}: 방향족 신호가 2H doublet 두 개뿐입니다. \textbf{7.52(t, 2H)와 7.44(t, 1H)는 생성물에만} 나타나므로, 이 두 신호의 적분을 기준으로 잔류 출발물질의 비율을 추정할 수 있습니다.
\item \textbf{Biphenyl}(호모커플링): COOH가 없고 7.3--7.7 영역에만 신호가 나옵니다.
\item \textbf{Benzoic acid}(탈브로민화 부반응 가능성): 일치환 패턴과 COOH를 함께 가지지만, AA$'$BB$'$ doublet 쌍이 없습니다.
\item \textbf{용매 잔류 신호 (DMSO-$d_6$)}: DMSO 2.50, \ce{H2O} 3.33, EtOH 1.06 (t) / 3.44 (q) / 4.35 (OH)(Fulmer 등, 2010). 재결정 용매인 EtOH가 보이면 건조 부족의 증거입니다.
\end{itemize}

% =====================================================================
\section{왜 반응 바이알인가 (L1 pp.\,14--15)}
L1 슬라이드 15의 비교를 이 실험에 맞춰 보면 다음과 같습니다. 부피가 약 11 mL로 작고, 70\,\si{\celsius}에서 30분 동안 \textbf{뚜껑을 닫아 증발 없이} 가열해야 합니다. 또 여러 학생의 바이알을 \textbf{공용 히팅 블록}에 동시에 꽂아야 합니다. 이런 조건에는 바이알이 적합합니다. RBF는 환류, 회전증발, Schlenk line처럼 진공이나 유리 접합부가 필요한 작업에 유리합니다. 실험 4(Grignard)에서 2구 RBF를 쓰는 것과 대비해 보세요.

% =====================================================================
\section{L1 연습문제: pimodivir의 11과 14}

\begin{center}
\setchemfig{atom sep=1.7em}
\begin{tabular}{c@{\hspace{16mm}}c}
\chemfig{*6(-N=-(-F)=-*5(-(-Bpin)=-N(-Ts)-)=)} &
\chemfig{*6(=N-(-Cl)=N-(-NH-R)=(-F)-)} \\[1mm]
\textbf{11} & \textbf{14} \quad {\footnotesize R = (2S,3S)-2-(ethoxycarbonyl)bicyclo[2.2.2]oct-3-yl}
\end{tabular}
\end{center}
\begin{itemize}[leftmargin=5mm,itemsep=0.5mm,topsep=0.5mm]
\item \textbf{11}: 5-fluoro-3-(pinacolatoboryl)-1-tosyl-7-azaindole. 브로마이드 \textbf{10}의 C3--Br이 Miyaura 붕소화(\ce{B2pin2}, \ce{Pd(OAc)2}/\ce{PPh3}, KOAc)로 C3--Bpin이 된 것입니다. \textbf{약염기 KOAc}를 쓰는 이유는, 생성된 Bpin이 남은 \textbf{10}과 Suzuki 커플링을 일으키는 것을 억제하기 위해서입니다(Ishiyama 등, 1995).
\item \textbf{14}: ethyl (2S,3S)-3-[(2-chloro-5-fluoropyrimidin-4-yl)amino]bicyclo[2.2.2]octane-2-carboxylate. 완제품 pimodivir는 azaindole C3과 pyrimidine C2가 연결된 구조입니다. 따라서 C2에 Cl이 있어야 합니다. 2-클로로피리미딘은 전자가 부족해 일반 아릴 클로라이드보다는 산화적 첨가가 쉬운 편이지만, 여전히 C--Br보다 느리고 기질에 질소가 많아 Pd에 달라붙을 수 있습니다. 그래서 전자가 풍부하고 부피가 큰 X-Phos 같은 리간드를 씁니다.
\item \textbf{15}에서 pimodivir로 가려면 ethyl ester가 COOH로, N-Ts가 N--H로 바뀌어야 합니다. 슬라이드의 마지막 두 조작(HCl, 이어 NaOH 수용액)이 이 변환을 담당하며, NaOH가 ester 가수분해를 맡는 것은 분명합니다. 구조는 Wikipedia/PubChem의 IUPAC명으로 대조했습니다. 각 단계 시약의 정확한 역할은 강의 해설로 확인하세요 \confirm{강의 해설}.
\end{itemize}

% =====================================================================
\section{보고서 고찰에 쓸 만한 질문}
\begin{enumerate}[leftmargin=6mm,itemsep=0.5mm]
\item 수율 손실은 어디서 생겼는가? 재결정 에탄올 사용량, 세척량, 여과지 잔류 등을 datasheet 기록과 연결해 논의합니다.
\item NMR로 본 순도: 잔류 4-bromobenzoic acid, biphenyl, EtOH/\ce{H2O}가 있는가? 적분비로 정량할 수 있는가?
\item TON/TOF 값의 의미와 한계(4절의 세 가지 가정).
\item \ce{PhB(OH)2}를 1.2 eq 쓰는 이유: protodeboronation과 호모커플링으로 일부가 소모되기 때문입니다.
\item 이 계는 얼마나 ``green''한가? 용매, 촉매량(148 ppm), 추출 생략이 장점이고, 산성화로 생기는 폐염과 Pd 폐액이 한계입니다.
\end{enumerate}

% =====================================================================
\section*{참고문헌}
{\footnotesize
\begin{enumerate}[leftmargin=6mm,itemsep=0.3mm,label={[\arabic*]}]
\item Hoang, T.\,G. \emph{CM3191 Chemical Synthesis Experiments, Organic Lab Manual}, AY26/27, Exp.\,1, pp.\,2--7; \emph{CM3191 Org L1: Suzuki--Miyaura Cross-Coupling in Water} (lecture slides).
\item Hamilton, A.\,E.; Buxton, A.\,M.; Peeples, C.\,J.; Chalker, J.\,M. An Operationally Simple Aqueous Suzuki--Miyaura Cross-Coupling Reaction for an Undergraduate Organic Chemistry Laboratory. \emph{J. Chem. Educ.} \textbf{2013}, \emph{90}, 1509--1513. doi:10.1021/ed4002333.
\item Chalker, J.\,M.; Wood, C.\,S.\,C.; Davis, B.\,G. A Convenient Catalyst for Aqueous and Protein Suzuki--Miyaura Cross-Coupling. \emph{J. Am. Chem. Soc.} \textbf{2009}, \emph{131}, 16346--16347. doi:10.1021/ja907150m.
\item Carrow, B.\,P.; Hartwig, J.\,F. Distinguishing Between Pathways for Transmetalation in Suzuki--Miyaura Reactions. \emph{J. Am. Chem. Soc.} \textbf{2011}, \emph{133}, 2116--2119. doi:10.1021/ja1108326.
\item Lennox, A.\,J.\,J.; Lloyd-Jones, G.\,C. Selection of Boron Reagents for Suzuki--Miyaura Coupling. \emph{Chem. Soc. Rev.} \textbf{2014}, \emph{43}, 412--443.
\item Miyaura, N.; Suzuki, A. Palladium-Catalyzed Cross-Coupling Reactions of Organoboron Compounds. \emph{Chem. Rev.} \textbf{1995}, \emph{95}, 2457--2483.
\item Ishiyama, T.; Murata, M.; Miyaura, N. Palladium(0)-Catalyzed Cross-Coupling Reaction of Alkoxydiboron with Haloarenes: A Direct Procedure for Arylboronic Esters. \emph{J. Org. Chem.} \textbf{1995}, \emph{60}, 7508--7510.
\item Roughley, S.\,D.; Jordan, A.\,M. The Medicinal Chemist's Toolbox. \emph{J. Med. Chem.} \textbf{2011}, \emph{54}, 3451--3479.
\item Fulmer, G.\,R. \emph{et al.} NMR Chemical Shifts of Trace Impurities. \emph{Organometallics} \textbf{2010}, \emph{29}, 2176--2179.
\item Pavia, D.\,L.; Lampman, G.\,M.; Kriz, G.\,S.; Vyvyan, J.\,R. \emph{Introduction to Spectroscopy}, Cengage (IR 작용기 흡수 범위).
\item Clark, M.\,P. \emph{et al.} Discovery of a Novel, First-in-Class, Orally Bioavailable Azaindole Inhibitor (VX-787) of Influenza PB2. \emph{J. Med. Chem.} \textbf{2014}, \emph{57}, 6668--6678. \quad 구조 확인: Pimodivir, PubChem CID 67286591.
\end{enumerate}
}

\end{document}
